\element{vitesses}{
  \begin{question}{vphi}\bareme{1}
    La vitesse de phase est définie comme
    \begin{multicols}{4}
      \begin{reponses}
        \bonne{$\frac{\omega}{\Re{\underline{k}}}$}
        \mauvaise{$\frac{1}{\Im{\underline{k}}}$}
        \mauvaise{$\frac{\omega}{\underline{k}}$}
        \mauvaise{$\frac{d\omega}{d\Re{\underline{k}}}$}
      \end{reponses}
    \end{multicols}
  \end{question}
}
\element{vitesses}{
  \begin{question}{vg}\bareme{1}
    La vitesse de groupe est définie comme
    \begin{multicols}{4}
      \begin{reponses}
        \mauvaise{$\frac{\omega}{\Re{\underline{k}}}$}
        \mauvaise{$\frac{1}{\Im{\underline{k}}}$}
        \mauvaise{$\frac{\omega}{\underline{k}}$}
        \bonne{$\frac{d\omega}{d\Re{\underline{k}}}$}
      \end{reponses}
    \end{multicols}
  \end{question}
}
\element{vitesses}{
  \begin{question}{vg}\bareme{1}
    La profondeur de peau est définie comme
    \begin{multicols}{4}
      \begin{reponses}
        \mauvaise{$\frac{\omega}{\Re{\underline{k}}}$}
        \bonne{$\frac{1}{\Im{\underline{k}}}$}
        \mauvaise{$\frac{\omega}{\underline{k}}$}
        \mauvaise{$\frac{d\omega}{d\Re{\underline{k}}}$}
      \end{reponses}
    \end{multicols}
  \end{question}
}
\setgroupmode{vitesses}{cyclic}

\element{equations}{
  \begin{questionmult}{dAlembert}\bareme{haut=2,p=0}
    Une équation de d'Alembert
    %\begin{multicols}{2}
      \begin{reponses}
        \mauvaise{s'écrit $\frac{\partial^2 y}{\partial x^2}+\frac{1}{c^2}\frac{\partial^2y}{\partial t^2}=0$}
        \mauvaise{est dispersive}
        \mauvaise{a pour relation de dispersion $\underline{k}^2=-j\frac{\omega}{D}$}
      \end{reponses}
    %\end{multicols}
  \end{questionmult}
}
\element{equations}{
  \begin{questionmult}{diffusion}\bareme{haut=2,p=0}
    Une équation de diffusion
    %\begin{multicols}{2}
      \begin{reponses}
        \mauvaise{s'écrit $\frac{\partial^2 y}{\partial t^2}+D\frac{\partial^2y}{\partial x^2}=0$}
        \bonne{est dispersive}
        \mauvaise{a pour relation de dispersion $k^2=\frac{\omega^2}{c^2}$}
      \end{reponses}
    %\end{multicols}
  \end{questionmult}
}

\element{dispersif}{
  \begin{questionmult}{dispersif}\bareme{haut=2,p=0}
    Dans un milieu dispersif,
    %\begin{multicols}{2}
      \begin{reponses}
        \bonne{Toutes les fréquences ne se propagnent pas à la même vitesse.}
        \bonne{Les paquets d'onde se déforment en se propageant.}
        \mauvaise{La relation de dispersion s'écrit $\omega = k c$ pour les ondes se propageant dans le sens des $x$ croissants.}
        \mauvaise{L'équation de propagation est forcément non-linéaire.}
      \end{reponses}
    %\end{multicols}
  \end{questionmult}
}